\section*{Bab \ref{ch:b3:canonical-ensemble} --- Ensembel Kanonik}

\begin{solution}{exo:b3:canonical-ensemble:1}
(a) $\mathcal P_2/\mathcal P_1 = (g_2/g_1)\eu^{-\beta\epsilon}$.
(b) $\beta\epsilon = 2.1/(8.62\times10^{-5} \times 2500) = 9.7$:
bagiannya $3\eu^{-9.7} \approx \num{2e-4}$. (c) Sebuah nyala membawa
$\sim\num{e15}$ atom natrium tiap cm$^3$: bahkan $10^{-4}$ di antaranya,
yang berdaur tiap beberapa nanodetik, mencurahkan $10^{19}$ foton kuning
tiap detik --- menyilaukan. (d) $3\eu^{-\beta\epsilon} = 0.1$:
$T = \epsilon/(k_{\text{B}}\ln30) \approx \qty{7200}{K}$ --- yakni
fotosfer bintang, bukan nyala api.
\end{solution}

\begin{solution}{exo:b3:canonical-ensemble:2}
(a) $z = 1 + \eu^{-\beta\epsilon} + \eu^{-2\beta\epsilon}$. (b)
$\langle E\rangle = \epsilon(\eu^{-\beta\epsilon} +
2\eu^{-2\beta\epsilon})/z$: $\to 0$ pada $T$ rendah, $\to \epsilon$
(yakni aras rata-ratanya) pada suhu tinggi. (c) $\mathcal P_1 = 1/(\eu^{\beta
\epsilon} + 1 + \eu^{-\beta\epsilon})$ tumbuh secara monoton bersama
$T$, dan mendekati batas atasnya $1/3$ hanya ketika $T \to \infty$. (d)
Bobot Boltzmann hanya menurun terhadap energi pada $T > 0$; pembalikan
populasi memerlukan \omterm{ex:b3:microcanonical-ensemble:negative}{suhu negatif}
\cref{ex:b3:microcanonical-ensemble:negative}, yang tak terjangkau oleh
tandon mana pun.
\end{solution}

\begin{solution}{exo:b3:canonical-ensemble:3}
(a) Faktor Boltzmann yang dikenakan pada $E_p = mgh$ pada $T$ seragam. (b) $h_0 =
k_{\text{B}}T/mg = \qty{8.4}{km}$. (c) $\eu^{-8848/8440} \approx
0.35$: sepertiga atmosfer --- karena itulah para pendaki puncak membawa
oksigen. (d) Atmosfer yang sebenarnya tidak isotermal: suhunya turun
seiring ketinggian, sehingga profilnya membelok menjauhi eksponensial
tunggal.
\end{solution}

\begin{solution}{exo:b3:canonical-ensemble:4}
(a) Tiga translasi: $\tfrac32 R = \qty{12.5}{J/(mol.K)}$, sesuai
persis. (b) Tambahkanlah dua rotasi: $\tfrac52 R =
\qty{20.8}{J/(mol.K)}$, sebagaimana terukur: getarannya beku. (c)
Ramalannya $\tfrac72 R = \qty{29.1}{J/(mol.K)}$; nilai terukur
$\approx26$ menunjukkan getarannya baru mencair sebagian ($\theta_{\text{vib}}
\approx \qty{3400}{K}$). (d) $3R = \qty{24.9}{J/(mol.K)}$:
Dulong--Petit, yang ditaati pada suhu ruang oleh sebagian besar logam
dan dilanggar oleh intan --- yakni kisah pembekuan
\cref{exo:b3:canonical-ensemble:6}.
\end{solution}

\begin{solution}{exo:b3:canonical-ensemble:5}
(a) $z = \eu^{-\beta\hbar\omega/2}/(1 - \eu^{-\beta\hbar\omega})$.
(b) $\langle E\rangle = -\partial_\beta\ln z$ memberikan bentuk yang
dinyatakan itu dengan $\langle n\rangle = 1/(\eu^{\beta\hbar\omega} - 1)$. (c)
$C \to k_{\text{B}}$ pada $T$ tinggi; $C \approx k_{\text{B}}(\beta
\hbar\omega)^2\eu^{-\beta\hbar\omega} \to 0$ pada suhu rendah. (d) Dengan
$\hbar\omega$ sebagai energi satu kuantum cahaya, $\langle n\rangle$
adalah cacah foton termal tiap ragamnya: hukum Planck tinggal satu bab
lagi.
\end{solution}

\begin{solution}{exo:b3:canonical-ensemble:6}
(a) Secara bilangan, $C = \tfrac12 \times 3Nk_{\text{B}}$ di dekat
$T/\theta_{\text{E}} \approx 0.34$. (b) $\theta_{\text{E}}/T =
4.3$: $C/3Nk_{\text{B}} = 4.3^2\eu^{-4.3}/(1 - \eu^{-4.3})^2
\approx 0.26$ --- intan pada suhu ruang baru memiliki seperempat
kapasitas kalor klasiknya: ``kejanggalannya'' adalah pembekuan kuantum
yang terpampang jelas. (c) $\theta_{\text{E}} = \qty{90}{K}$
menempatkan timbal jauh di dalam rezim klasik pada \qty{300}{K}. (d)
Anggapan frekuensi tunggalnya: zat padat yang sebenarnya memiliki
spektrum ragam sampai ke bunyi berpanjang gelombang besar, yang
kuantanya murah lalu memberikan hukum $T^3$ --- yakni perbaikan Debye,
dalam \cref{ch:b3:photons-phonons}.
\end{solution}

\begin{solution}{exo:b3:canonical-ensemble:7}
(a) $f(v) \propto v^2\eu^{-mv^2/2k_{\text{B}}T}$: $v_{\text{p}} =
\sqrt{2k_{\text{B}}T/m}$: \qty{413}{m/s} untuk N$_2$,
\qty{1540}{m/s} untuk H$_2$. (b) $mv_{\text{esc}}^2/2k_{\text{B}}T
\approx 730$ untuk N$_2$, $52$ untuk H$_2$. (c) $\eu^{-730}$ berarti
tak pernah; $\eu^{-52} \sim \num{e-23}$ tiap waktu tinggalnya lambat ---
tetapi atmosfer atasnya yang panas ($\sim\qty{1000}{K}$) melunakkan
eksponen hidrogennya menjadi $\sim15$: hidrogen bocor sepanjang waktu
geologi, nitrogen bertahan. (d) Laju lepas dan suhu eksosfernya: sumur
\qty{60}{km/s} Jupiter membuat eksponen hidrogen pun menjadi
astronomis --- raksasa gas menyimpan apa yang dilepaskan dunia kecil
yang hangat.
\end{solution}

\begin{solution}{exo:b3:canonical-ensemble:8}
(a) $\partial_\beta^2\ln Z = \langle E^2\rangle - \langle
E\rangle^2$, dan $C = \partial_T\langle E\rangle =
-k_{\text{B}}\beta^2\partial_\beta\langle E\rangle$. (b) Kedua
ruasnya memberikan $Nk_{\text{B}}(\beta\epsilon)^2\eu^{\beta\epsilon}/(
\eu^{\beta\epsilon} + 1)^2$. (c) $\langle E\rangle \propto N$,
$\Delta E \propto \sqrt N$. (d) Seberapa kuat sebuah sistem menanggapi
pemanasan sama dengan seberapa besar energinya bergetar secara spontan
--- tanggapan dan fluktuasi adalah dua bacaan atas turunan kedua yang
sama, yakni pola (fluktuasi--disipasi) yang berulang di seluruh fisika.
\end{solution}

\begin{solution}{exo:b3:canonical-ensemble:9}
(a) $\ln2 = E_a\,\Delta T/k_{\text{B}}T^2$ dengan $\Delta T =
\qty{10}{K}$, $T = \qty{298}{K}$: $E_a =
0.693\,k_{\text{B}}T^2/10 \approx \qty{0.55}{eV}$. (b) Dari
$298$ ke \qty{278}{K}: faktor $\eu^{E_a(1/278 - 1/298)/k_{\text{B}}}
\approx 4.4$ lebih lambat --- karena itulah lemari es mengawetkan
makanan. (c) $373 \to \qty{366}{K}$: lajunya turun sebesar
$\approx1.4$: telur sebelas menit memerlukan seperempat jam. (d) Reaksi
dimenangkan oleh ekor eksponensial molekul di atas penghalangnya:
geserlah suhunya sedikit dan populasi ekornya bergeser luar biasa ---
nilai rata-ratanya hampir tak berarti.
\end{solution}

\begin{solution}{exo:b3:canonical-ensemble:10}
(a) $f_{\text{terlipat}} = 1/(1 + \eu^{-\beta(\Delta E -
T\Delta S)})$ dengan entropi keadaan terbukanya sudah terlipat ke dalam
energi bebasnya. (b) Pada $T_{\text{m}} = \Delta E/\Delta S$ kedua
energi bebasnya seri: separuh dan separuh. (c) $T_{\text{m}} =
\num{4.8e-19}/\num{1.38e-21} = \qty{348}{K}$ (\qty{75}{\celsius});
lebarnya $\delta T \sim k_{\text{B}}T_{\text{m}}^2/\Delta E \approx
\qty{3.5}{K}$. (d) Kekooperatifan mengalikan $\Delta E$ maupun
$\Delta S$ dengan cacah sentuhan yang putus bersamaan, sehingga
$T_{\text{m}}$ tetap tetapi lebarnya menyusut $\propto 1/\Delta E$:
untai DNA terpisah dalam rentang beberapa kelvin, dan itulah yang
membuat mesin PCR dapat berdaur secara bersih antara ``melebur'' dan
``menempel''.
\end{solution}

\begin{solution}{exo:b3:canonical-ensemble:11}
(a) $z = 2\cosh(\beta\mu B)$; $M = N\mu\tanh(\beta\mu B) \approx
N\mu^2B/k_{\text{B}}T$ untuk argumen kecil: yakni $1/T$ Curie. (b)
$\mu_{\text{B}}B/k_{\text{B}}T$: $\num{2.2e-3}$ pada \qty{300}{K};
$0.67$ pada \qty{1}{K} --- dari acuh tak acuh menjadi tersearahkan
kuat. (c) $S$ hanya bergantung pada $\mu B/k_{\text{B}}T$; menurunkan
$B$ pada entropi tetap menyeret $T$ turun secara sebanding: $\qty{1}{K} \to
\qty{0.1}{K}$. (d) Ketika $k_{\text{B}}T$ mencapai energi antarspinnya,
entropinya tak lagi dikendalikan medan:
$\mu_0\mu_{\text{B}}^2/4\pi a^3 \approx \qty{7e-26}{J} \sim
\qty{5}{mK}$ --- yakni lantai klasiknya (momen inti, seribu kali lebih
lemah, mendorongnya sampai ke mikrokelvin).
\end{solution}

\begin{solution}{exo:b3:canonical-ensemble:12}
(a) $z_{\text{rot}} = \sum_J(2J+1)\eu^{-\beta BJ(J+1)}
\to \int(2J+1)\eu^{-\beta BJ(J+1)}\dd J = k_{\text{B}}T/B$: maka
$\langle E\rangle = k_{\text{B}}T$ dan $C_{\text{rot}} =
k_{\text{B}}$. (b) H$_2$: $\theta_{\text{rot}} = \qty{87}{K}$ ---
anak tangganya terletak dalam jangkauan laboratorium; N$_2$: \qty{2.9}{K},
yang baru membeku di dekat helium cair, sehingga nitrogen selalu
menunjukkan $\tfrac52 R$. (c) Dataran $\tfrac32 R$ (di bawah $\sim\qty{50}{K}$),
$\tfrac52 R$ (dari $\sim200$ sampai $\sim\qty{1000}{K}$), lalu menanjak
menuju $\tfrac72 R$ di dekat $\theta_{\text{vib}} \approx
\qty{6000}{K}$. (d) $J$ ganjil dan genap termasuk jenis spin inti yang
berbeda, yang saling berubah secara lambat, sehingga $C_V$ hidrogen
dingin bergantung pada riwayat orto--paranya --- statistik inti yang
diaudit oleh sebuah kalorimeter.
\end{solution}

\begin{solution}{pb:b3:canonical-ensemble:1}
\textbf{1.} $V = \tfrac43\pi r^3 = \qty{4.0e-20}{m^3}$: $m' = V
\Delta\rho = \qty{8.3e-18}{kg}$, beratnya $m'g = \qty{8.1e-17}{N}$.
\textbf{2.} $n(h) = n_0\,\eu^{-m'gh/k_{\text{B}}T}$: yakni hukum
barometrik, yang disusutkan sampai muat pada kaca objek mikroskop.
\textbf{3.} $h_0 = k_{\text{B}}T/m'g = \num{4.04e-21}/\num{8.1e-17}
\approx \qty{50}{\micro m}$.
\textbf{4.} $h_0 \propto 1/r^3$: sebaran jejari sebesar $10\%$ menjadi
sebaran tinggi skala sebesar $30\%$ --- butiran yang beragam ukurannya
mengaburkan tangganya menjadi bubur. Perrin memilahnya berbulan-bulan
dengan sentrifugasi berulang.
\textbf{5.} Rumus yang sama, dengan massa berselisih $10^{10}$:
``atmosfer'' butirannya $10^{10}$ kali lebih dangkal --- kilometer
menyusut menjadi puluhan mikrometer, dan justru itulah yang membuatnya
teramati utuh di bawah mikroskop.
\textbf{6.} Nisbah berturutannya $0.55$, $0.55$, $0.57$: tetap dalam
ralat pencacahannya --- eksponensial. $h_0 =
\qty{30}{\micro m}/\ln(100/55) \approx \qty{50}{\micro m}$.
\textbf{7.} $k_{\text{B}} = m'g\,h_0/T = \num{8.1e-17} \times
\num{5.0e-5}/293 \approx \qty{1.4e-23}{J/K}$.
\textbf{8.} $N_{\text{A}} = R/k_{\text{B}} \approx
\num{6.0e23}\,\unit{mol^{-1}}$.
\textbf{9.} Di sini melesetnya beberapa persen saja (dengan data yang
diidealkan); percobaan Perrin yang sebenarnya berserak sebesar $\pm15\%$
di sekitar nilai modernnya --- luar biasa untuk butiran yang dicacah
dengan tangan, dan sepenuhnya menentukan bagi orde besarnya.
\textbf{10.} $m_{\text{H}} = \qty{e-3}{kg/mol}/N_{\text{A}} =
\qty{1.7e-27}{kg}$.
\textbf{11.} Kimia makroskopis menyediakan $R$; mikroskopi cahaya dan
penimbangan menyediakan $m'$; pencacahan menyediakan tangganya: tiga
pengukuran di atas meja bersama-sama menentukan massa sebuah atom yang
tak dapat dilihat siapa pun.
\textbf{12.} Tak ada: penurunannya hanya memakai ``sistem yang bertukar
energi dengan sebuah tandon'' --- faktor Boltzmann buta terhadap ukuran,
dan persis itulah yang diperiksa Perrin.
\textbf{13.} $D = k_{\text{B}}T/6\pi\eta r =
\num{4.04e-21}/\num{4.0e-9} \approx \qty{e-12}{m^2/s}$.
\textbf{14.} $\sqrt{2Dt} = \sqrt{\num{1.2e-10}} \approx
\qty{11}{\micro m}$ tiap menitnya: dapat diukur dengan nyaman memakai
kisi lensa mata dan kesabaran.
\textbf{15.} Dua gejala yang saling bebas, dua rumus yang saling bebas,
satu bilangan: kesepakatan antara tangga statis dan getaran dinamis tak
menyisakan celah bagi kebetulan --- hipotesis molekulernya meramalkan
keduanya.
\textbf{16.} $\sqrt{k_{\text{B}}T/m} = \sqrt{\num{4.04e-21}/
\num{4.8e-17}} \approx \qty{9}{mm/s}$ tiap komponennya.
\textbf{17.} Butirannya dihantam $10^{19}$ kali tiap detik lalu
melupakan kecepatannya dalam jarak nanometer: penerbangan balistiknya
tak teramati, dan yang dilihat mata adalah integralnya --- yakni jalan
acak yang mendifusi.
\textbf{18.} Seretan Stokes (kekentalan fluida jilid Tahun 2)
menyediakan $6\pi\eta r$-nya; \omterm{thm:b3:canonical-ensemble:boltzmann}{ensembel kanoniknya} menyediakan
$k_{\text{B}}T$-nya: $D$ Einstein adalah hasil baginya, yakni mekanika
dan statistika dalam satu pecahan.
\textbf{19.} Khayalan tak dapat dicacah: begitu $N_{\text{A}}$ menjadi
nisbah dua bilangan terukur, lengkap dengan palang ralatnya, atom
menjadi obyek percobaan --- Ostwald mengaku kalah secara tertulis pada
1909.
\textbf{20.} Hamburan biru langit, elektrolisis dengan muatan elektron
terukur, helium yang tertimbun dari radium: empat saluran fisis yang tak
berkaitan memusat pada satu $6\times10^{23}$, sehingga bilangan itu
menjadi sifat alam, bukan sifat teori mana pun.
\textbf{21.} Energi gravitasi yang menarik butirannya ke bawah, entropi
susunan yang menyebarkannya ke atas, dipertukarkan pada nilai tukar $T$:
tangganya \emph{adalah} minimum $F = E - TS$.
\textbf{22.} $h_0 \propto 1/g$: pada $10^5g$,
$h_0 = k_{\text{B}}T/(m'\times10^5g) = \num{4.04e-21}/
\num{9.8e-17} \approx \qty{40}{\micro m}$ untuk proteinnya --- yakni
kesetimbangan pengendapan, percobaan Perrin yang dijalankan setiap hari
di jurusan biokimia.
\textbf{23.} Dari ``terlalu tinggi'' ($h_0 \gg$ kedalaman medannya:
tanpa gradien yang kasat mata) sampai ``terlalu tipis'' ($h_0 \lesssim r$):
jejari yang berguna terbentang kira-kira $0.1$--$\qty{0.5}{\micro m}$
--- pilihan Perrin bukan keberuntungan melainkan rancangan.
\textbf{24.} Dengan $k_{\text{B}}$ yang kini eksak menurut definisi,
percobaan yang sama mengalibrasi butirannya (ukuran atau massa jenisnya)
atau mengaudit perangkatnya: fisikanya --- tangga Boltzmann --- tak
tersentuh; yang berpindah hanyalah besaran mana yang dianggap tak
diketahui.
\textbf{25.} Populasi butiran \qty{0.2}{\micro m} berkurang separuh tiap
$\sim\qty{35}{\micro m}$; dibaca lewat $n_0\eu^{-m'gh/k_{\text{B}}
T}$, tangganya menghasilkan $N_{\text{A}} \approx \num{6e23}$: bilangan
Avogadro yang dicacah butir demi butir --- dan perdebatan atomnya
ditutup di atas meja mikroskop.
\end{solution}
